\setcounter{section}{-1}
\Section{Символика}

Применение символики \textit{не регулируется строго}, её можно писать и произносить
по-разному.  

Множества обозначают заглавными латинскими буквами (\(A\), \(B\), \(C\), \(\dots\), \(Z\)),
элементы множеств~--- \textit{прописными латинскими буквами}
(\(a\), \(b\), \(c\), \(\dots\), \(z\)).

Иногда для обозначения числовых множеств будут использоваться прописные греческие
буквы (\(\alpha\), \(\beta\), \(\gamma\), \(\dots\), \(\omega\)).

Приведём примеры некоторых обозначений.

\vspace{-8pt}
\begin{table}[H]
\begin{tabular}{ll}
\textit{Обозначение} & \textit{Смысл}\\
\(a \in A\) & <<\(a\) является элементом множества \(A\)>> или <<\(a\) принадлежит множеству
    \(A\)>>.\\
\(a \notin A\) & <<\(a\) не является элементом множества \(A\)>> или <<\(a\) не принадлежит
    множеству \(A\)>>\\
\(A \cup B\) & <<объединение \(A\) и \(B\)>>\\
\(A \cap B\) & <<пересечение \(A\) и \(B\)>>\\
\(A \setminus B\) & <<разность \(A\) и \(B\)>>\\
\(A \subset B\) & <<\(A\) является подмножеством множества \(B\)>>\\
\(A \not\subset B\) & <<\(A\) не является подмножеством множества \(B\)>>\\
\(\forall\) & квантор всеобщности (<<для любого>>, <<для всех>>)\\
\(\exists\) & квантор существования (<<существует>>)\\
\(\Rightarrow\) & импликация (<<следует>>, <<то>>, <<тогда>>)\\
\(\Leftrightarrow\) & эквивалентность (<<тогда и только тогда>>, <<необходимо и достаточно>>)\\
\(\lor\) & дизъюнкция (<<или>>)\\
\(\&\) & конъюнкция (<<и>>)\\
\(\neg\) & отрицание (<<не>>)\\
\(\N\), \(\Z\), \(\Q\), \(\R\) & числовые множества \\
\(\eqdef\) & <<по определению>> \\
\(\mapsto\) & <<выполняется>> \\
\(\colon\) & <<такой, что>> \\
\end{tabular}
\end{table}
\vspace{-14pt}

\begin{example}
Приведем пример использования символики:
\[
    \left[\N \in \Z\right] \& \left[\Z \in \Q\right] \Rightarrow \left[\N \in \Q\right]
\]
\end{example}

\Subsection{Примеры, использование символики}

\begin{example}
<<множество \(A\) является подмножеством множества \(B\)>> по определению то же самое, что
и <<для любого \(x\), принадлежащего множеству \(A\), выполняется, что \(x\) принадлежит
множеству \(B\)>>:
\[
    \left[A \subset B\right] \eqdef \left[\forall x \in A \mapsto x \in B\right]   
\]
\end{example}

\begin{example}
<<множество \(A\) равняется множеству \(B\)>> по определению то же самое, что и
<<множество \(A\) является подмножеством множества \(B\) \textbf{и}
множество \(B\) является подмножеством множества \(A\)>>:
\[
    \left[A = B\right] \eqdef \left[A \subset B\right] \& \left[B \subset A\right]
\]
На практике, теоретико-множественные равенства доказываются обычно следующим образом:
\[
    \left[A = B\right] = \left[\forall x \in A \mapsto x \in B\right] \&
    \left[\forall x \in B \mapsto x \in A\right]
\]
\end{example}

\begin{example}
<<Объединением можеств \(A\) и \(B\)>> называют множество, все элементы \(x\) которого
принадлежат множеству \(A\) \textbf{или} множеству \(B\):
\[
    \left[A \cup B\right] = \left\{x \colon \left[x \in A\right] \lor \left[x \in B\right]\right\}
\]
\end{example}

\begin{example}
<<Пересечением можеств \(A\) и \(B\)>> называют множество, все элементы \(x\) которого
принадлежат множеству \(A\) \textbf{и} множеству \(B\):
\[
    \left[A \cap B\right] = \left\{x \colon \left[x \in A\right] \& \left[x \in B\right]\right\}
\]
\end{example}

\begin{example}
<<Разностью можеств \(A\) и \(B\)>> называют множество, все элементы \(x\) которого
принадлежат множеству \(A\) \textbf{и не} принадлежат множеству \(B\):
\[
    \left[A \setminus B\right] = \left\{x \colon \left[x \in A\right] \&
    \left[x \not\in B\right]\right\}
\]
\end{example}

\begin{example}
При отрицании, квантор всеобщности~\(\forall\) заменяется на квантор существования~\(\exists\). Так
же после квантора всеобщности~\(\forall\) принято ставить~<<\(\mapsto\)>>, а после квантора
существования~\(\exists\)~--- <<\(\colon\)>> (это правило не является строгим, символы стоит использовать в
зависимости от их произношения).

Обозначим \(\mathfrak{A} = \left[A \subset B\right] = \left[\forall a \in A \mapsto a \in B
\right]\), тогда:
\[
    \neg \mathfrak{A} = \left[A \not\subset B\right] =
    \left[\exists a \in A\colon a \notin B\right]
\]
\end{example}

\begin{example}
При отрицании, все символы, которые находятся при кванторах~\(\forall\), \(\exists\) не меняются,
а все остальные~--- \textit{заменяются на противоположные} (\(\in\)~--- на \(\notin\), \(>\)~--- на
\(\leq\), \(\lor\)~--- на \(\&\) и~т.\,д.).

Обозначим \(\mathfrak{B} = \left[\forall x \in X \mapsto \lvert x \rvert < \alpha \right]\), тогда:
\[
    \neg \mathfrak{B} = \left[\exists x \in X\colon \lvert x \rvert \ge \alpha\right]
\]
\end{example}

\begin{example}
Теперь рассмотрим пример с несколькими кванторами

Обозначим \(\mathfrak{C} = \left[\exists \alpha > 0\colon \forall x \in X \mapsto \lvert x \rvert <
\alpha\right]\), тогда:
\[
    \neg \mathfrak{C} = \left[ \forall \alpha > 0 \; \exists x_\alpha \in X\colon
    \lvert x_\alpha \rvert \geq \alpha\right]
\]

Стоит обратить внимание на то, что при отрицании мы добавили к~\(x\) индекс~\(\alpha\). Это связано
с тем, что для \textit{каждого~\(\alpha\) будет существовать свой~\(x_{\alpha}\)}, так что
использование~\(x\) в данном случае является некорректным.
\end{example}

Стоит так же отметить тот факт, что запись через символы является довольно свободной, и у неё нет
строгих правил записи.


\Section{Теория действительного числа}
\Subsection{Рациональные числа и их свойства}

\begin{definition}
\textit{Рациональным числом} называется дробь~\(\dfrac{p}{q}\), где \(q \in \N\), \(p \in \Z\).
\end{definition}

Запишем основные свойства рациональных чисел:

\begin{properties}
\item Если \(a = b\) и \(b = c\), то \(a = c\),

Если \(a > b\) и \(b > c\), то \(a > c\).

(\textit{транзитивность <<\(=\)>> и <<\(>\)>>})

\item \(a + b = b + a\) (\textit{коммутативность сложения});
\item \((a + b) + c = a + (b + c)\) (\textit{ассоциотивность сложения});
\item Существует число~\(0\) такое, что для любого~\(a\) выполняется \(a + 0 = a\) (\textit{особая
    роль нуля});
\item Для любого числа~\(a\) найдется противоположное число~\(-a\) такое, что \(a + (-a) = 0\)
    (\textit{существование противоположных чисел}).
\item \(a \cdot b = b \cdot a\) (\textit{коммутативность умножения});
\item \((a \cdot b) \cdot c = a \cdot (b \cdot c)\) (\textit{ассоциативность умножения});
\item Существует число 1 такое, что для любого~\(a\) выполняется \(a \cdot 1 = a\) (\textit{особая
    роль 1});
\item Для любого~\(a \neq 0\) существует обратное число \(\dfrac{1}{a}\) такое, что \(a \cdot
    \dfrac{1}{a} = 1\) (\textit{существование обратных чисел});
\item \((a + b) \cdot c = a \cdot c + b \cdot c\) (\textit{дистрибутивность сложения относительно
    умножения});
\item Если \(a > b\), то для любого~\(c\) выполняется \(a + c > b + c\);
\item Если \(a > b\), то для любого~\(c > 0\) выполняется \(a \cdot c > b \cdot c\);
\item Для любого~\(a > 0\) найдется \(n \in \mathbb{N}\) такое, что \(n > a\) (\textit{аксиома
    Архимеда}).
\end{properties}

\begin{mtopic}{Плотность множества рациональных чисел}
Если \(a > b\), то найдется \(c \in \mathbb{Q}\) такое, что \(a > c > b\).

Для доказательства этого утверждения достаточно взять \(c = \dfrac{a + b}{2}\).
\end{mtopic}

\Section{Сечения и иррациональные числа}

\begin{definition}
\textit{Сечением} множества \(\Q\) называется разбиение \(\Q\) на два непустых множества
\(A^*\) и \(A_*\) таких, что:
\begin{properties}
\item \(A_* \cup A^* = \Q\)
\item \(A_* \cap A^* = \emptyset\)
\item \(\left[\forall x \in A_*\right] \& \left[\forall y \in A^*\right] \mapsto
    y > x\)
\end{properties}
\end{definition}

\begin{example} Приведем пример сечения:

\(A_* =  \left\{a\colon \frac{1}{10^3} > a\right\}\),
\(A^* = \left\{a\colon a \geq \frac{1}{10^3}\right\}\)
\end{example}

\textit{Сечение} обозначается как \(\left(A_*, A^*\right)\), где \(A_*\)~---
\textit{нижний класс сечения}, а \(A^*\)~--- \textit{верхний класс сечения}.

\begin{example}
\(b, c \in \N\colon (b + 1)^2 > c > b^2\) (\(c\) не является квадратом какого-либо
натурального числа):
\begin{gather}
A_* = \left\{a\colon \left[a \leq 0\right] \lor \left[c > a^2\right]\right\} \\
A^* = \left\{a\colon \left[a > 0\right] \& \left[a^2 \geq c\right]\right\}
\end{gather}
\end{example}

\begin{proposition}
Сечение \(\left(A_*, A^*\right)\) примера 11 таково, что \textit{нижний} класс сечения
не имеет \textit{наибольшего} числа, а \textit{верхний} класс сечения не
имеет \textit{наибольшего} числа.
\end{proposition}

\begin{proposition}
\textit{Не существует} \(\left(A_*, A^*\right)\) такого, что \(A_*\) имеет
\textit{наибольшее} число, а \(A^*\) имеет \textit{наименьшее} число
(так как нарушится или первое, или второе условие из определения сечения).
\end{proposition}

\textbf{Рассмотрим сечения множества рациональных чисел, которые существуют:}

\begin{properties}
\item В \(A_*\) \textbf{нет} наибольшего числа, в \(A^*\) \textbf{есть} наименьшее число
\item В \(A_*\) \textbf{есть} наибольшее число, в \(A^*\) \textbf{нет} наименьшего числа
\item В \(A_*\) \textbf{нет} наибольшего числа, в \(A^*\) \textbf{нет} наименьшего числа
\end{properties}

Так как сечения \(1^\circ\) и \(2^\circ\) определяются одним и тем же рациональным числом,
то будем использовать сечение \(1^\circ\).

\begin{definition}
\textit{Иррациональным числом} назовем сечение вида \(3^\circ\).
\end{definition}

Множество иррациональных чисел обозначается как \(\J\),
а рациональные числа будем записывать через греческие буквы \(\alpha = \left(A_*, A^*\right)\).

\begin{definition}
\textit{Действительным числом} называется любое сечение \(\Q\) вида \(1^\circ\) или~\(2^\circ\).
\end{definition}

Действительные числа будем обозначать греческими буквами \(\alpha =
\left(A_*, A^*\right)\), \(\R = \Q \cup \J\)

\textbf{Важно отметить}, что если известно \(A_*\), то известно и всё сечение
\(\left(A_*, A^*\right)\) 

Пусть \(\alpha = \left(A_*, A^*\right)\), \(\beta = \left(B_*, B^*\right)\).

\begin{definition}
\(
    \left[\alpha = \beta\right] \eqdef \left[A_* = B_*\right]
\)
\end{definition}

\begin{definition}
\(
    \left[\alpha \not = \beta\right] \eqdef \left[A_* \not = B_*\right]
\)
\end{definition}

\begin{proposition}
\(
    \left[\alpha \not = \beta\right] \Rightarrow \left[A_* \subset B_*\right] \lor
    \left[B_* \subset A_*\right]
\)
\end{proposition}

\begin{definition}
\(
    \left[\alpha > \beta\right] \eqdef \left[B_* \subset A_*\right] \&
    \left[B_* \not = A_*\right]
\)
\end{definition}